\section{Conditional confinement in Bernoulli coordinates}
\label{v9:confinement:section}

Two different statements will be proved. Low noise gives half-amplitude
confinement on the acquired probability range $[0.40,0.53]$ by one
continuous polynomial argument. A separate signed-deletion calculation
covers the full acquired noise range on $[0.41,0.42]$. Neither statement
asserts the subsequent strict count comparison. Both preserve the
orientation selected for the fixed-target maximizing pair.

\subsection{Conditional statement and acquired application}
Retain the fixed-target hypotheses of Section~\ref{v9:resources:section}:
$n\ge42$, nonnegative excess $a$, $c=\CE$, closed-support
contact, and the true predecessor $C_{n-1}\le c$. The actual raw threshold
$H$ can have either sign. Assume an actual central gap $(-A,B)$ with
$A,B>0$, and write
\[
 p=\Pp(X>0),\quad q=1-p,\quad m=\E X_+=\E X_-,\quad
 \bar A=m/q,\quad \bar B=m/p,\quad \ell=\bar A+\bar B,
\]
\[
 X=\ell(Y-p),\qquad Y=B_0+U,\qquad
 B_0=\mathbf1_{\{X>0\}},\qquad \E(U\mid B_0)=0.
\]
Define $U_i$ as the conditional value of $U$ on $\{B_0=i\}$.
Assume the already acquired bounds
\begin{equation}\label{v9:confinement:entry}
 p\in[2/5,3/5],\qquad |U|<3/5,\qquad
 d=1-(A+B)/\ell<9/125,
 \qquad \lambda=n\E U^2\le3/40.
\end{equation}
The gap implies $A\le\bar A$, $B\le\bar B$, and hence $d\ge0$.

\begin{theorem}\label{v9:confinement:theorem}
Under these hypotheses, the entire conditional supports satisfy
$|U|<1/2$. The conclusion is uniform in all raw $H$ and all $n\ge42$.
It includes the closed noise endpoint $\lambda=3/40$ and zero noise.
\end{theorem}

The probability interval in this abstract theorem is $[0.4,0.6]$.
The currently acquired entry supplies the incoming strip and inward bounds
only on the union of panels 40--52, namely $[0.40,0.53]$. The application
below is restricted to that smaller interval. In particular this section
does not acquire the strip on $(0.53,0.6]$ or treat the larger-noise region.

\subsection{One common moment enclosure}
Put $v=pq$, $s=\E U^2=\lambda/n$, and
$\theta=\E(X-\E[X\mid B_0])^2$. Conditional centering and variance
give the exact identities
\begin{equation}\label{v9:confinement:identities}
 \ell^2=(v+s)^{-1},\qquad
 \theta=\frac{s}{v+s}=\frac{\lambda}{nv+\lambda},\qquad
 \ell^2=\frac{1-\theta}{v}.
\end{equation}
Throughout \eqref{v9:confinement:entry}, $6/25\le v\le1/4$ and
$\max(p,q)\le3/5$. Use the rational constants
\[
 \Theta=\frac{37}{5000},\qquad E=\frac{249}{125},\qquad
 r=\frac{221}{200}.
\]
Then
\[
 \theta\le\frac{5}{677}<\Theta,\qquad
 \ell^2\ge4(1-\Theta)>E^2,
\]
with exact differences
$\Theta-5/677=49/3385000$ and $4(1-\Theta)-E^2=73/31250$.

Because $B_0$ is the actual sign of $X$, the absolute cubic
expansion has no sign approximation:
\begin{equation}\label{v9:confinement:cubic}
 \rho:=\E|X|^3=\ell^3\left\{v(1-2v)
       +3\E[|B_0-p|U^2]
       +\E[\operatorname{sgn}(B_0-p)U^3]\right\}.
\end{equation}
The linear term vanishes by conditional centering. The incoming strip,
used before the conclusion, gives $|U|^3\le(3/5)U^2$ and
$|B_0-p|\le3/5$. Also $\ell^3\le v^{-3/2}$. Therefore
\[
 \rho\le\rho_B(v)+\frac{12/5}{v^{3/2}}\frac\lambda n,
 \qquad \rho_B(v)=\frac{1-2v}{\sqrt v}.
\]
The function $\rho_B$ decreases on $(0,1/4]$, since its derivative is
$-\tfrac12v^{-3/2}-v^{-1/2}<0$. The exact squared comparison
\[
 24^2(6/25)^3-(12/5)^2=\frac{34416}{15625}>0
\]
gives $\rho\le\rho_B(6/25)+24\lambda/n<r$. For the last step,
\[
 r-\frac3{70}=\frac{1487}{1400}>0,\qquad
 \left(r-\frac3{70}\right)^2\frac6{25}
       -\left(\frac{13}{25}\right)^2=\frac{8707}{24500000}>0.
\]

If $\theta_\pm$ are the two nonnegative conditional variance contributions
in physical coordinates, their sum is $\theta$. Expanding the signed
second moment gives
\[
 \beta:=\E[X|X|]=(1-\theta)(1-2p)+\theta_+-\theta_-.
\]
Consequently, for the actual $p$, including $p>1/2$,
\begin{equation}\label{v9:confinement:beta}
 1-2p-2(1-p)\Theta\le\beta\le1-2p+2p\Theta.
\end{equation}
All these bounds concern one common marginal. No assertion about
independent conditional laws or a moment relaxation is used.

\subsection{All-threshold contact and exterior rays}
\begin{lemma}[Uniform heat bound and ray continuation]\label{v9:confinement:heat}
Under fixed-target contact, nonnegative excess and the sharp predecessor,
let $r\ge\rho$ and let $b$ bound the directional signed moment: $\beta$
on the positive ray, $-\beta$ on the negative ray. At a radius $u>\sqrt3$,
the polynomials \eqref{v9:confinement:Q}--\eqref{v9:confinement:P} below
satisfy the following implication. If $Q_{r,b}(u)<0$ and
$P_{r,b}(u)<0$, that boundary is absent from the closed support. If also
$Q',Q'',L,L',L''<0$ there, with the moment cap held fixed, the entire
closed exterior ray is absent. This assertion requires no probability
panel or incoming noise strip and is uniform in all raw $H$ and $n\ge42$.
\end{lemma}
\begin{proof}
For $w=n-1$, $s_n=\sqrt{n/w}$, $F_w(t)=\Phi(t/\sqrt w)$ and
$e_n=n(s_n-1)-1/2$, fixed-target contact and the predecessor bound give,
at every closed-support point $t$,
\begin{equation}\label{v9:confinement:contact}
 0\le\mathcal H_n(H,t)
       +c\{-|t|^3+\tfrac32\rho t^2+3\beta t+\rho e_n\}
       -na,
\end{equation}
where
\[
 \frac{\mathcal H_n(H,t)}{n^{3/2}}
 =F_w(H-t)-F_n(H)+tF_n'(H)-\frac{t^2-1}{2}F_n''(H).
\]
The term $-na$ is nonpositive and may be dropped. This is a
fixed-target statement; no equality between $c$ and the actual score is
assumed.

The spatial and variance expansions were derived in
\eqref{v9:entry:heat}, together with each of their four remainder bounds.
They apply for every $u\ge1$, before any conditional strip is used.
For $t=\pm u$, write $w=n-1$, $z=H/\sqrt w$,
$A_u=(u^3-3u)/6$ and $H_4(u)=u^4-6u^2+3$. Then
\[
 \frac{\mathcal H_n(H,\pm u)}{s_n^3}
 =\pm A_u(1-z^2)\phi(z)
       +\frac{H_4(u)}{24\sqrt w}\phi'''(z)+R,
\]
\[
 |R|\le\frac{u^5/100+3u/20}{w}
      +\frac{51(3u^2-2)}{960w^{3/2}}.
\]
For $u>\sqrt3$, $A_u>0$. The bounds
$\|(1-z^2)\phi(z)\|_\infty=\phi(0)$ and
$\|\phi'''\|_\infty\le2/\pi<2\phi(0)$ control both signs.
Unlike the initial large-radius calculation, we retain $|H_4(u)|$;
its sign is not assumed positive.

Set $S=81/80$, $c_0=819/2000<c$ and $e=8/875$.
For $n\ge42$, $s_n<S$, $\sqrt w>6$, and $e_n<e$.
For completeness,
$ne_n=s_n^2(2s_n+1)/(2(s_n+1)^2)$ is increasing for $s_n\ge1$;
its value at $S$ is $793881/2073680<48/125$.
Define
\begin{align}
 E_u&=u^5/4100+3u/820+17(3u^2-2)/78720,\notag\\
 Q_{r,b}(u)&=-u^3+\tfrac32ru^2+3bu+er,\label{v9:confinement:Q}\\
 P_{r,b}(u)&=\tfrac25S^3\{A_u+|H_4(u)|/72\}
                  +S^3 E_u+c_0Q_{r,b}(u).\label{v9:confinement:P}
\end{align}
If the directional signed moment is at most $b$, $Q<0$ licenses replacing
the positive multiplier $c$ by $c_0$. Then $P<0$ contradicts
\eqref{v9:confinement:contact} at that boundary for every raw $H$.

To continue along a physical ray, hold $r,b$ fixed. The derivative
estimate established in the proof of Theorem~\ref{v9:entry:initial-support} is
\[
 |\partial_t\mathcal H_n|<\tfrac25(t^2/2+1/2+e).
\]
Put $L_{r,b}(u)=\tfrac25(u^2/2+1/2+e)+c_0Q'_{r,b}(u)$.
If $Q,Q',Q'',L,L',L''$ are negative at a boundary, the signs persist
along the whole ray: $Q'''=-6$ and $L''=2/5-6c_0<0$.
Thus the contact upper bound strictly decreases there. Boundary atoms
and the entire closed exterior ray are excluded.
\end{proof}

\subsection{A continuous one-parameter certificate}
For the negative ray use $z=p$ and the signed moment $-\beta$; for the
positive ray use $z=1-p$ and $\beta$. In both cases $z\in[2/5,3/5]$,
and \eqref{v9:confinement:beta} gives the directional cap
\[
 b_z=2z-1+2(1-z)\Theta.
\]
This folds the actual probability, not the endpoints of a probability
panel. It does not reflect the chosen discrepancy or change its law.
Take the boundary $u=E(z+1/2)$. Thus
\[
 u\in[u_0,u_1]=[2241/1250,2739/1250],\qquad
 b_z=\alpha u+\delta,
 \quad\alpha=\frac{4963}{4980},\quad\delta=-\frac{9889}{5000}.
\]
Here $u_0^2>3$, so $H_4'(u)=4u(u^2-3)>0$;
$H_4(u_1)<-2.75$ proves $H_4<0$ on this whole interval.
Let $\mathcal P(u)=P_{r,\alpha u+\delta}(u)$.
Replacing $|H_4|$ by $-H_4$ in \eqref{v9:confinement:P} gives an explicit quintic
$\mathcal P(u)=\sum_{j=0}^5 a_j u^j$, with
\begin{align*}
 a_5&=S^3/4100,& a_4&=-\tfrac25S^3/72,\\
 a_3&=\tfrac25S^3/6-c_0,&
 a_2&=\tfrac25S^3/12+51S^3/78720+c_0(3r/2+3\alpha),\\
 a_1&=-\tfrac25S^3/2+3S^3/820+3c_0\delta,&
 a_0&=-\tfrac25S^3/24-34S^3/78720+c_0er.
\end{align*}
The exact rational endpoint inequalities are
\[
 \mathcal P'''(u_0)<-2.24,\qquad \mathcal P'''(u_1)<-2.27,
 \qquad \mathcal P'(u_0)>0.915,\quad \mathcal P'(u_1)>0.745.
\]
Indeed $\mathcal P'''$ is a convex quadratic, with second derivative
$120a_5>0$. Its negative endpoint values make it negative throughout.
Therefore $\mathcal P'$ is concave and exceeds the chord between its
positive endpoint values. It follows that $\mathcal P$ increases and
\begin{equation}\label{v9:confinement:margin}
 \mathcal P(u)\le\mathcal P(u_1)
 =-\frac{4560620039489147528734839}{25625000000000000000000000}
 <-0.1779.
\end{equation}
One stronger-looking shortcut is false:
$\mathcal P''(u_0)>0.0228$. It is not used.

The multiplier and ray signs require a separate calculation.
The boundary family
\[
 \mathcal Q(u)=Q_{r,\alpha u+\delta}(u)
 =-u^3+(3r/2+3\alpha)u^2+3\delta u+er
\]
has total derivatives $\mathcal Q''(u_0)<-1.46$ and
$\mathcal Q'(u_1)=179223/6250000>0.0286$.
Since $\mathcal Q''$ decreases, $\mathcal Q'$ decreases and stays
positive. Hence $\mathcal Q$ increases, and
$\mathcal Q(u)\le\mathcal Q(u_1)<-1.198$.

For each physical ray the directional cap is fixed. Its derivatives,
evaluated across the boundary family, are instead
\begin{align*}
 q_1(u)&=-3u^2+(3r+3\alpha)u+3\delta,&
 q_2(u)&=-6u+3r,\\
 l(u)&=\tfrac25(u^2/2+1/2+e)+c_0q_1(u),&
 l_1(u)&=\tfrac25u+c_0q_2(u),\qquad l_2=\tfrac25-6c_0.
\end{align*}
They are not the total derivatives of $\mathcal Q$.
The following rational upper bounds certify their signs on the interval:
\begin{center}
\begin{tabular}{lrr}\toprule
family & value at $u_0$ & total derivative at $u_0$\\\midrule
$q_1$ & $<-4.272$ & $<-4.452$\\
$q_2$ & $<-7.441$ & $=-6$\\
$l$   & $<-0.903$ & $<-1.105$\\
$l_1$ & $<-2.330$ & $=-2.057$\\\bottomrule
\end{tabular}
\end{center}
The two quadratic families have negative constant second derivatives
$q_1''=-6$ and $l''=-2.057$; the others are affine decreasing.
Thus all four stay negative, as does $l_2=-2.057$.
Together with $\mathcal Q<0$, these are the six required signs at every
boundary. Holding its actual directional cap fixed, the preceding ray
argument excludes the whole exterior ray beyond each boundary.

\subsection{Completion, application and remaining premises}
For each actual $p$ the two excluded rays start at
\[
 u_-=E(p+1/2),\qquad u_+=E(3/2-p).
\]
Because $\ell>E$, the remaining support satisfies
$-1/2<Y=X/\ell+p<3/2$ strictly. The already proved actual gap supplies
\[
 U_0\le(\bar A-A)/\ell\le d<9/125,
 \qquad U_1\ge-(\bar B-B)/\ell\ge-d>-9/125.
\]
For $B_0=0$ the new lower bound is $U_0>-1/2$;
for $B_0=1$ the new upper bound is $U_1<1/2$.
This proves Theorem~\ref{v9:confinement:theorem}. If $\lambda=0$, then $\E U^2=0$
and $U=0$ directly. There is no division by noise or limiting argument.

\begin{corollary}\label{v9:confinement:application}
Retain the accepted fixed-target entry and the ordered elementary
moment/gap contraction on its 30 probability panels. On its acquired
subrange $p\in[0.40,0.53]$, all actual $n\ge42$ with
$\lambda\in[0,3/40]$ satisfy $|U|<1/2$ for every raw $H$.
The previously acquired inward bound $d<0.072<1/4$ is retained.
\end{corollary}
\begin{proof}
That antecedent provides the strip $|U|<0.6$ on the 13 closed panels
40--52 and $d\le0.0719539227<0.072$.
Their closed union is $[0.40,0.53]$. Apply the theorem at the actual $p$.
\end{proof}

\subsection{The central panel at larger noise}
\label{v9:confinement:high-setup}
We now restrict to the oriented panel $p\in[l,h]=[0.41,0.42]$ and use
\eqref{v9:resources:central-entry}. Set
\[
 N=42,\quad\Theta=0.054,\quad d_0=0.069,\quad
 E=1.9706,\quad F_\ell=2.034,\quad \sigma_0=0.6,
\]
\[
 r_0=1/2,\quad\tau=1/5,\quad\zeta=10^{-5},\quad
 \kappa=997/15000<\phi(0)/6,
\]
\[
 c_l=0.4095<c=\CE<c_u=0.40974.
\]
Here $E$ is a lower scale bound and $\sigma_0$ is the incoming strip
width, not $\sqrt{10}$. These constants replace the older floor-47,
width-$0.55$ confinement parameters; every gate below uses floor 42
and the already acquired width $0.6$.

Use distinct count-coordinate inner endpoints
$a_0=p-A/\ell$, $b_0=p+B/\ell$.
They are support points with $0\le a_0\le d$, $1-d\le b_0\le1$ and
$b_0-a_0=1-d$. The inner inequalities $U_0\le a_0$, $U_1\ge b_0-1$
are weak; strictness comes from $d<d_0$. Let
\[
 x=np+H/\ell,\quad \nu=\operatorname{Law}\left(x-\sum_{i<n}Y_i\right),
 \qquad J_{\rm gap}=\int_{[a_0,b_0)}(y-a_0)\nu(dy).
\]
The central gap contributes $m\ell J_{\rm gap}=(1-\theta)J_{\rm gap}$
to the scalar leakage, so
\begin{equation}\label{v9:confinement:leakage}
 \mathcal L\ge(1-\theta)J_{\rm gap}.
\end{equation}
Write $g(y)=\nu([y,\infty))$ and let $Q_c$ be the count-coordinate
contact majorant, with $g\le Q_c$ everywhere and equality on the closed
support. Differentiating fixed-target contact gives
\begin{equation}\label{v9:confinement:contact-derivative}
 \sqrt n(-Q_c'(y))=\ell\phi(H/\sqrt n)
 +\frac\ell n\{3c\beta+(H\phi(H/\sqrt n)+3c\rho)t-3ct|t|\},
 \qquad t=\ell(y-p).
\end{equation}
The excess changes only an additive constant. The multiplier here is
$c$, not an upper bound on the actual score.

\subsubsection{One resource for displacement, noise and leakage}
\label{v9:confinement:simplex}
Define the nonnegative endpoint distances
\[
 M_+=\E[X-B;X>0]=m-pB,\qquad M_-=\E[-X-A;X<0]=m-qA.
\]
With $\sigma=\sqrt{10}$, the endpoint expansion of the actual defect gives
\begin{align*}
 D\ge{}&m\{\sqrt{\sigma+1}A-\sqrt{\sigma-1}B\}^2
 +2(\sigma+2)B^2M_++2(\sigma+4)A^2M_-\\
 &+3(\sigma+1)B\theta_++3(\sigma+3)A\theta_-.
\end{align*}
Indeed put $Z=X-B\ge0$ on the positive class. Then
\[
 T_+=2B^2M_++3B\E[Z^2;X>0]+\E[Z^3;X>0],\quad
 R_+=BM_++\E[Z^2;X>0].
\]
The combination $(\sigma-1)T_++6BR_+$ gives the stated positive
terms because $\E[Z^2;X>0]\ge\theta_+$. The negative expansion is the
same magnitude calculation with $\sigma+1$. The remaining constant
term is the square since $2\sqrt{(\sigma+1)(\sigma-1)}=6$.
Discard this nonnegative square and use $A\ge E(l-d_0)$,
$B\ge E(1-h-d_0)$. The positive coefficients
\begin{align*}
 c_+&=2(5.16)E^2(1-h-d_0)^2,&c_-&=2(7.16)E^2(l-d_0)^2,\\
 b_*&=\min\{3(4.16)E(1-h-d_0),3(6.16)E(l-d_0)\}
\end{align*}
satisfy
\begin{equation}\label{v9:confinement:endpoint-cost}
 D\ge c_+M_++c_-M_-+b_*\theta,\qquad
 d=\frac{M_+/p+M_-/q}{\ell},\qquad
 \epsilon:=\E|U|\le\frac{2(M_++M_-)}\ell.
\end{equation}
The last bound follows signwise from
$\E|Z-\E Z|\le2\E Z$ for $Z\ge0$.

On a closed noise band $[\lambda_0,\lambda_1]$, let $S_\varepsilon$
be the whole-rectangle resource from \eqref{v9:resources:rectangle}.
Because $n\theta=\ell^2\lambda\ge E^2\lambda_0$, the single common
budget implies
\begin{equation}\label{v9:confinement:resource-simplex}
 nc_+M_++nc_-M_-+n^{3/2}\mathcal L/\kappa\le R_*,
 \qquad R_*=S_\varepsilon/\kappa-b_*E^2\lambda_0.
\end{equation}
A strictly negative $R_*$ excludes a law; equality is retained.
With $\alpha=\min\{c_+lE,c_-(1-h)E\}>0$ and $G=R_*/\alpha$, the
same simplex gives $nd\le G$. These costs are simultaneous, so the
defect budget cannot be spent independently on each one.

\subsubsection{Concentration at both deleted indices}
\label{v9:confinement:concentration}
Independently couple the two conditional noises and sample an independent
Bernoulli label. This realizes the actual marginal as
$Y=U_0+B_0(1+U_1-U_0)$, with increment at least $1-d>0$.
Conditional on all pairs, an interval of diameter strictly less than
$1-d$ selects an antichain of label sets. For an antichain $\mathcal A$
in $j$ labels, a uniformly random permutation gives a chain meeting it
at most once. Taking expectations yields
\[
 \sum_{A\in\mathcal A}\binom j{|A|}^{-1}\le1,\qquad
 \Pp(\mathcal A)=\sum_{A\in\mathcal A}\binom j{|A|}^{-1}
        \binom j{|A|}p^{|A|}q^{j-|A|}
 \le\max_k\binom jk p^kq^{j-k}.
\]
Conditioning back proves the bound for the actual iid sum, also after
an independent shift. Strict diameter is essential: a closed interval
equal to a Bernoulli increment can contain both atoms.

Both actual deleted indices satisfy $j\ge n-2\ge40$. Put
\[
 r_L=l-(1-l)/40,\quad r_R=h+h/40,\qquad
 v_*=\min\{r_L(1-r_L),r_R(1-r_R)\}.
\]
The gate $0<r_L<r_R<1$ holds. Adjacent binomial-mass ratios give
$(j+1)p-1\le k\le(j+1)p$ for a mode; hence $r=k/j\in[r_L,r_R]$
and the mode is interior. The classical finite Stirling inequalities
$\sqrt{2\pi j}(j/e)^j<j!<
\sqrt{2\pi j}(j/e)^j e^{1/(12j)}$ give
\[
 \sqrt j\binom jk p^kq^{j-k}
 <\frac{\phi(0)e^{1/(12j)-jD(r\Vert p)}}{\sqrt{r(1-r)}}
 <\frac{0.4}{\sqrt{v_*}(1-1/480)}<M.
\]
Here entropy is nonnegative by Jensen's inequality, and
$e^z\le(1-z)^{-1}$ for $0\le z<1$ follows termwise from its power
series. Choose $M$ by taking the downward $10^{-10}$ grid square root of
$0.16/[v_*(1-1/480)^2]$ and adding one grid unit; thus $M$ is strictly
larger than the exact root.
Thus every interval of strict width below $1-d$ has probability less
than $M/\sqrt j$. No bound on $C_{n-2}$ is assumed. Finally
\begin{equation}\label{v9:confinement:deleted-factor}
 F_d=1.037,\qquad F_d^2>\max\{(42/41)^3,42/40\}.
\end{equation}
These ratios decrease with $n$ and control the distance factor from
$n-1$ summands and the noise/pair-event factor from $n-2$ summands.

\subsubsection{Signed deletion, including endpoint atoms}
\label{v9:confinement:deletion}
Let $m_n=n-1$, $Z=x-\sum_{i=1}^{m_n}Y_i$,
$\mu=x-m_np=p+H/\ell$, and let $K$ count the labels. For a removed label
define $T_i(I)=\Pp(x-\sum_{k<m_n}Y_k-i-U_i\in I)$.
Exchangeability and $m_np-K=Z-\mu+\sum U_k$ give exactly
\begin{align*}
 \nu(I)-T_1(I)&=\frac{\E[(Z-\mu+\sum U_k)\ind_I(Z)]}{m_np},\\
 T_0(I)-\nu(I)&=\frac{\E[(Z-\mu+\sum U_k)\ind_I(Z)]}{m_nq}.
\end{align*}
For a removed full pair, the remaining conditional interval probability
is $f(Y_k)\in[0,M/\sqrt{n-2}]$. Since $\E U_k=0$,
\[
 \E[U_kf(Y_k)]\le\frac{M}{\sqrt{n-2}}\E(U_k)_+
                  =\frac{\epsilon M}{2\sqrt{n-2}}.
\]
The signed term gives the factor $\epsilon/2$ and also covers
deterministic conditional noise.

For independently coupled noises let $\Delta=U_1-U_0$.
Then $\Delta\ge-d$, $\E\Delta=0$, and
$\operatorname{Var}\Delta\le s/p$, since $p<q$ and $s=\lambda/n$.
On $\Delta\le\tau$, the point for $T_1$ in $I$ maps to the point for
$T_0$ in $I^+=I+1+[-d,\tau]$. Conditional concentration, Cantelli and
Markov applied to $\Delta+d\ge0$ give
\begin{align}
 T_1(I)&\le T_0(I^+)+\frac M{\sqrt{n-2}}\Pp(\Delta>\tau),\notag\\
 \Pp(\Delta>\tau)&\le\min\left\{\frac{s}{s+p\tau^2},\frac d{d+\tau}\right\}
 \le B_\tau:=\min\left\{\frac{\lambda_1}{\lambda_1+Nl\tau^2},
                         \frac G{N\tau+G}\right\}.
 \label{v9:confinement:bad}
\end{align}
The exceptional event is strict; $\Delta=\tau$ is in the good event.
All intervals below have strict widths below $1-d$. If $D_I,D_J$
bound the positive parts of their upper endpoints minus $\mu$, the two
deletion identities imply
\begin{equation}\label{v9:confinement:deletion-error}
 \sqrt n\,\nu(I)\le\sqrt n\,\nu(I^+)+A_*/n+
             a_e\epsilon+B_*^{\rm err}\Pp(\Delta>\tau),
\end{equation}
where
\[
 a_e=\frac{F_dM}{2}\left(\frac1l+\frac1{1-h}\right),\quad
 A_*=F_dM\left(\frac{D_I}l+\frac{D_J}{1-h}\right),\quad
 B_*^{\rm err}=F_dM.
\]
Indeed $\sqrt n/(n-1)^{3/2}\le F_d/n$ and
$\sqrt{n/(n-2)}\le F_d$. Half-open intervals may be enlarged to closed
ones for upper concentration bounds; no CDF continuity is assumed.

\subsubsection{Whole-band slope bounds}
\label{v9:confinement:slopes}
Fix $|H|\in[h_0,h_1]$ with its original sign and a noise band
$[\lambda_0,\lambda_1]$. Let $r_B$ be the outward upper bound for
$(1-2l(1-l))/\sqrt{l(1-l)}$ obtained by dividing its numerator by the downward $10^{-10}$ grid
square root of $l(1-l)$, and set
\[
 \rho_*=\min\{1.22,r_B+24\lambda_1/N\},\quad
 \beta_-=0.09736,\quad\beta_+=0.234,
\]
\[
 \varphi_-=0.39886(1-h_1^2/(2N))>0,\qquad
 v_H^+=\begin{cases}0.4h_1,&H\ge0,\\-\varphi_-h_0,&H\le0.\end{cases}
\]
The coefficient 24 is licensed by the incoming $0.6$ strip:
$24^2[l(1-l)]^3>[3(1-l)+0.6]^2$. It follows from the same-law
cubic expansion before the smaller strip is acquired. The density
floor follows from $e^{-z}\ge1-z$; $v_H^+$ bounds $H\phi(H/\sqrt n)$
in the correct direction for both signs.

Take $B^*=1.2>\sqrt{(1-l)/l}$ and $f=0.4F_\ell$. Define
\begin{align*}
 C_{{\rm sl},-}&=F_\ell\pos{(v_H^++3c_u\rho_*)^2/(12c_l)-3c_l\beta_-},\\
 C_{{\rm sl},+}&=F_\ell\pos{0.4h_1B^*+3c_uB^*(B^*-1)-3c_l\beta_-},\\
 L_g&=E\varphi_--\max(C_{{\rm sl},-},C_{{\rm sl},+})/N,\\
 U_+&=f+\frac{F_\ell}{N}
       \left\{3c_u\beta_++\frac{(0.4h_1+3c_u\rho_*)^2}{12c_l}\right\}.
\end{align*}
For $t\le0$, the bracket in \eqref{v9:confinement:contact-derivative}
is bounded below by
$3c_l\beta_-+(v_H^++3c_u\rho_*)t+3c_lt^2$.
The gate $v_H^++3c_u\rho_*>0$ permits this square completion.
On the central positive interval $0\le t\le B^*$ use $\rho\ge1$ and
$t(t-1)\le B^*(B^*-1)$ for $t\ge1$; for $t\le1$ the latter term
has favorable sign. Thus $L_g$ is a lower slope throughout the central
gap. On the full positive half-line, completing the downward quadratic
gives the upper slope $U_+$. Negative corrections in lower bounds and
positive corrections in upper bounds use $\ell/n\le F_\ell/N$.

The lower exterior interval has a sharper slope. Put
\[
 t_*=-E(l+\tau+\zeta),\quad B(t)=\beta_-+\rho_*t+t^2,\quad
 B_{\min}=B(\min\{t_*,-\rho_*/2\}).
\]
For $t\le t_*<0$,
\[
 3c(\beta+\rho t+t^2)+H\phi t
 \ge3c_lB(t)+v_H^+t+3(c_u-c_l)\min(0,B_{\min}).
\]
First use $\beta\ge\beta_-$, $\rho\le\rho_*$ and $t<0$; when varying
$c$, only the possible negative part of $B$ incurs the final correction.
Minimize the remaining upward quadratic at
$\min\{t_*,-(v_H^++3c_l\rho_*)/(6c_l)\}$ and denote the total minimum
by $q_*$. Then
\[
 L=E\varphi_--F_\ell\pos{-q_*}/N
\]
is a lower slope on the entire lower interval. The finite gates require
$L>0$, $L_g>0$ and $U_+>0$ on every box.

\subsubsection{Two strict interval contradictions}
\label{v9:confinement:intervals}
For a putative support point $a_0-r$, $r\ge r_0$, take
\[
 I=[a_0-r,a_0-d-\tau-\zeta],\qquad
 I^+=[a_0-r+1-d,b_0-\zeta].
\]
The incoming strip gives $r<\sigma_0+d_0$. The exact gates
$g_0=1-\sigma_0-2d_0>0$ and $\sigma_0+d_0<1-d_0$ place
$I^+\subset[a_0+g_0,b_0)$ and make both widths strictly less than $1-d$.
The intervals are disjoint. Contact at the left endpoint and the
global majorant at the right endpoint give
\[
 \sqrt n\,\nu(I)\ge L(r-d-\tau-\zeta),\qquad
 \nu(I^+)\le J_{\rm gap}/g_0.
\]
Indeed
$\nu([a_0-r,a_0-d-\tau-\zeta))\ge
Q_c(a_0-r)-Q_c(a_0-d-\tau-\zeta)$; a closed interval only increases
its mass. No contact at the right endpoint is needed.
For the directed distances use
\[
 v_H=\begin{cases}-h_0/F_\ell,&H\ge0,\\h_1/E,&H\le0,\end{cases}\qquad
 D_I=\pos{v_H-l-\tau-\zeta},\quad D_J=\pos{v_H+1-l-\zeta}.
\]
These are independent of $r$. Substitution in the deletion bound shows
that a lower support point would require
\begin{equation}\label{v9:confinement:lower-necessary}
 L(r-\tau-\zeta)\le Ld+a_e\epsilon+\sqrt nJ_{\rm gap}/g_0+
                     A_*/n+B_*^{\rm err}\Pp(\Delta>\tau).
\end{equation}
As $L>0$, a strict contradiction at $r=r_0$ excludes the entire lower
part of the acquired strip.

For a putative support point $y\ge1+r_0$, take
\[
 I=[a_0,a_0+\tau],\qquad I^+=[b_0,a_0+1+2\tau].
\]
The strict gates $d_0+2\tau<r_0$, $1-d_0-\tau>0$ and
$1-2d_0-2\tau>0$ give disjointness, widths below $1-d$, and
$I^+\subset[b_0,y)$. Since
$\nu([a_0,b_0))=Q_c(a_0)-Q_c(b_0)$ and
$\nu([a_0+\tau,b_0))\le J_{\rm gap}/\tau$, contact implies
\[
 \sqrt n\,\nu(I)\ge L_g(1-d)-\sqrt nJ_{\rm gap}/\tau,\qquad
 \sqrt n\,\nu(I^+)\le U_+(\sigma_0+d).
\]
The last bound uses contact at $b_0,y$, includes an atom at $b_0$ and
excludes the contacting atom at $y$. The strict containment is essential.
For this side use
\[
 D_I=\pos{d_0+\tau+h_1/E-l},\qquad
 D_J=\pos{1+d_0+2\tau+h_1/E-l}.
\]
An upper support point would therefore require
\begin{equation}\label{v9:confinement:upper-necessary}
 L_g-\sigma_0U_+\le(L_g+U_+)d+a_e\epsilon+
            \sqrt nJ_{\rm gap}/\tau+A_*/n+B_*^{\rm err}\Pp(\Delta>\tau).
\end{equation}

\subsubsection{Two error bounds on the same resource simplex}
\label{v9:confinement:errors}
For either side its right-hand error is
$a_dd+a_e\epsilon+a_J\sqrt nJ_{\rm gap}+A_*/n+
B_*^{\rm err}\Pp(\Delta>\tau)$, with
\[
 (a_d,a_J)=
 \begin{cases}(L,1/g_0),&\text{lower side},\\
              (L_g+U_+,1/\tau),&\text{upper side}.\end{cases}
\]
Equations \eqref{v9:confinement:endpoint-cost} and
\eqref{v9:confinement:leakage} give the three linear costs per unit of
\eqref{v9:confinement:resource-simplex}:
\[
 z_+=\frac{a_d/l+2a_e}{Ec_+},\qquad
 z_-=\frac{a_d/(1-h)+2a_e}{Ec_-},\qquad z_J=\frac{\kappa a_J}{1-\Theta}.
\]
For example the positive contribution is at most
$(a_d/l+2a_e)M_+/E=z_+(nc_+M_+)/n$; the leakage contribution is at
most $z_J(n^{3/2}\mathcal L/\kappa)/n$. Let $z=\max(z_+,z_-,z_J)$.
The gate $zR_*+A_*>0$ licenses replacing $n$ by $N$, giving
\begin{equation}\label{v9:confinement:error0}
 \mathcal E_0=(zR_*+A_*)/N+B_*^{\rm err}B_\tau.
\end{equation}
This includes retained $R_*=0$ boxes; resource equality is not excluded.

A second bound pays for the pair event inside the same simplex. At the
three nonzero vertices, upper linear-cost/displacement pairs are
\begin{align*}
 (v_+,w_+)&=\left(z_+R_*/N,\frac{R_*}{Ec_+l}\right),&
 (v_-,w_-)&=\left(z_-R_*/N,\frac{R_*}{Ec_-(1-h)}\right),\\
 (v_J,w_J)&=(z_JR_*/N,0).
\end{align*}
They bound the linear error and $nd$. The nonnegative origin contributes
zero and can be included without changing their maximum. For $t,x\ge0$,
\[
 \frac{x}{N\tau+x}\le
 \frac{N\tau}{(N\tau+t)^2}x+\frac{t^2}{(N\tau+t)^2},
\]
because the difference is $N\tau(x-t)^2/[(N\tau+x)(N\tau+t)^2]\ge0$.
Use $d/(d+\tau)=nd/(n\tau+nd)\le nd/(N\tau+nd)$. Every tangent gives
\begin{equation}\label{v9:confinement:errort}
 \mathcal E_t=\frac{A_*}{N}+
 \max_{i\in\{+,-,J\}}\left\{v_i+B_*^{\rm err}
                 \frac{N\tau}{(N\tau+t)^2}w_i\right\}
 +B_*^{\rm err}\frac{t^2}{(N\tau+t)^2}.
\end{equation}
Take the minimum of \eqref{v9:confinement:error0} and the eleven valid
bounds \eqref{v9:confinement:errort} at $t=jG/10$, $j=0,\ldots,10$.
Each bounds one actual resource simplex; no incompatible maximizing
allocations are combined.

\begin{condition}[Complete central-panel confinement gates]
\label{v9:confinement:high-condition}
Use both signs of $H$, eight closed phase bands
$[907j/4000,907(j+1)/4000]$, $j=0,\ldots,7$, and eight closed noise
bands $[3k/40,3(k+1)/40]$, $k=0,\ldots,7$.
Retain all 128 boxes in the coverage list. On each use the complete
incident-cell resource of Condition~\ref{v9:resources:phase-condition}.
Exclude only an empty incident list, strictly negative $S_\varepsilon$
or strictly negative $R_*$. Otherwise check all domains, directed roots,
positive slopes and inverse-index gates above, then require both strict
inequalities, with the error computed for the respective side,
\begin{equation}\label{v9:confinement:strict-gates}
 L(r_0-\tau-\zeta)>\mathcal E,\qquad
 L_g-\sigma_0U_+>\mathcal E.
\end{equation}
For the lowest noise band the all-threshold low-noise theorem may
replace these two ray gates.
\end{condition}

\begin{theorem}[Complete conditional confinement on the central panel]
\label{v9:confinement:central}
Under the fixed-target entry and contact hypotheses, the panel
contraction and signed-resource conditions, and
Condition~\ref{v9:confinement:high-condition}, every actual $n\ge42$ with
$p\in[0.41,0.42]$ satisfies $|U|<1/2$ on the entire conditional
supports and $d<0.069<1/4$.
\end{theorem}
\begin{proof}
The acquired $|H|<1.814$ and $\lambda<0.6$ lie inside the closed cover.
The 16 first-band boxes are covered by
Corollary~\ref{v9:confinement:application}, uniformly in $H$.
On the 112-box complement the exact gate evaluation has this census;
the finite inequalities themselves, not just the count, are required:
\begin{center}
\begin{tabular}{lr}\toprule
necessary exclusion or strict conclusion&boxes\\\midrule
no incident actual moment cell&14\\
negative whole-rectangle resource&5\\
noise lower-endpoint cost exceeds the resource&65\\
both strict contact-ray inequalities&28\\\bottomrule
\end{tabular}
\end{center}
All 56 inequalities in the last row have strict margin greater than
$0.00892011917$. The lower contradiction extends to every
$r\in[r_0,\sigma_0+d_0)$ because its error is independent of $r$ and
$L>0$. The upper contradiction already applies to every support point
$y\ge1+r_0$ in the acquired strip. No support point therefore lies at
or below $a_0-r_0$ or at or above $1+r_0$. Since $a_0\ge0$,
$b_0\ge1-d$ and $d<d_0<r_0$, both conditional noises have strict
amplitude below $r_0$. Boundary atoms are included in the contradictions.
The seam $3/40$ belongs to both adjacent closed bands; both phase signs
are retained at $H=0$.
\end{proof}

The finite gates in the final condition were checked by exact replay
and a separate reconstruction of the interval, slope and common-resource
formulas. This computational statement is separate from the analytic
derivation. The unassisted confinement diagnostics remain failed
sufficient bounds in their recorded low-noise boxes; the composition
uses the proved low-noise theorem there without changing those failures.

Thus low-noise confinement is proved conditionally on $[0.40,0.53]$,
and complete acquired-noise confinement conditionally on $[0.41,0.42]$.
Higher noise on the other panels and strict count comparison remain
separate obligations. In particular, count formulas requiring $p\le1/2$
are not transferred to above-half panels by reflecting the already
chosen discrepancy.
